\documentclass[10pt]{extarticle}
\usepackage[a4paper,margin=1.5cm]{geometry}
\usepackage[T1]{fontenc}
\usepackage{amsmath,amssymb}
\usepackage{graphicx}
\usepackage[export]{adjustbox}
\usepackage{multicol}
\usepackage{enumitem}
\usepackage{booktabs,array,tabularx}
\setlength{\columnsep}{0.7cm}
\setlength{\parindent}{0pt}
\setlength{\parskip}{3pt}
\pagestyle{empty}
\makeatletter
\newcommand{\qfigdef}[4]{\@namedef{qf@file@#1}{#2}\@namedef{qf@w@#1}{#3}\@namedef{qf@wide@#1}{#4}}
\newcommand{\qf@check}[1]{\@ifundefined{qf@file@#1}{\PackageError{score2}{Unknown figure id #1}{}}{}}
\newcommand{\qf@img}[4][]{\includegraphics[width=#2\dimexpr\@nameuse{qf@w@#4}pt\relax,max width=#3,max height=0.45\textheight,#1]{\@nameuse{qf@file@#4}}}
\newif\ifqf@full
\newif\ifqf@boxed
\newsavebox{\qf@box}
\newcommand{\qf@open}{\global\qf@fulltrue\par\global\qf@boxedtrue
  \setbox\qf@box\vbox\bgroup\hsize\textwidth\linewidth\textwidth}
\newcommand{\qf@emit}[1]{\par\egroup\global\qf@boxedfalse
  \ifdim\dimexpr\ht\qf@box+\dp\qf@box\relax<#1\textheight\noindent\usebox{\qf@box}\par\else\unvbox\qf@box\fi}
\newcommand{\qfullsplit}{\ifqf@boxed\qf@emit{0.85}\fi}
\newcommand{\qf@close}{\ifqf@boxed\qf@emit{0.45}\fi\global\qf@fullfalse}
\newcommand{\qfullbegin}{\end{multicols}\qf@open}
\newcommand{\qfullend}{\qf@close\begin{multicols}{2}}
\newcommand{\qfullstart}{\qf@open}
\newcommand{\qfulldone}{\qf@close}
\newcommand{\qfig}[2][1]{\qf@check{#2}%
  \ifqf@full\qf@blk{#1}{#2}\else
  \ifnum\@nameuse{qf@wide@#2}=1 \expandafter\@firstoftwo\else\expandafter\@secondoftwo\fi
  {\end{multicols}\begin{center}\qf@img{#1}{\textwidth}{#2}\end{center}\begin{multicols}{2}}%
  {\qf@blk{#1}{#2}}\fi}
\newcommand{\qf@blk}[2]{\par\smallskip{\centering\qf@img{#1}{\linewidth}{#2}\par}\smallskip}
\newcommand{\qfigi}[2][1]{\qf@check{#2}\qf@img[valign=c]{#1}{\linewidth}{#2}}
\makeatother
\newcommand{\papersec}[1]{\par\bigskip{\centering\Large\bfseries PAPER-#1\par}\medskip}
\newcommand{\partsec}[2]{\par\medskip\noindent{\large\bfseries PART-#1: #2}\par\smallskip}
\newcommand{\sect}[1]{\par\smallskip\noindent\textbf{#1}\par}
\newcommand{\qn}[1]{\par\smallskip\noindent\textbf{#1.}~}
\newcommand{\paraq}[2]{\par\smallskip\noindent\textbf{Paragraph for Questions #1 and #2}\par}
\newenvironment{qoptions}{\begin{itemize}[leftmargin=1.8em,itemsep=1pt,topsep=2pt,label={}]}{\end{itemize}}
\newcommand{\opt}[2]{\item[(#1)] #2}
% figure registry (generated by score2_tools.py finalize)
\qfigdef{p1_phy_s1_q01_1}{JEE Advanced (2021 Series) - Test 06 (22 Apr 2021)-img/p1_phy_s1_q01_1.png}{99.5}{0}
\qfigdef{p1_phy_s1_q02_1}{JEE Advanced (2021 Series) - Test 06 (22 Apr 2021)-img/p1_phy_s1_q02_1.png}{134.2}{0}
\qfigdef{p1_phy_s1_q02_2}{JEE Advanced (2021 Series) - Test 06 (22 Apr 2021)-img/p1_phy_s1_q02_2.png}{134.8}{0}
\qfigdef{p1_phy_s1_q03_1}{JEE Advanced (2021 Series) - Test 06 (22 Apr 2021)-img/p1_phy_s1_q03_1.png}{178.0}{0}
\qfigdef{p1_phy_s1_q04_1}{JEE Advanced (2021 Series) - Test 06 (22 Apr 2021)-img/p1_phy_s1_q04_1.png}{141.4}{0}
\qfigdef{p1_phy_s1_q06_1}{JEE Advanced (2021 Series) - Test 06 (22 Apr 2021)-img/p1_phy_s1_q06_1.png}{80.1}{0}
\qfigdef{p1_phy_s2_para7_8_1}{JEE Advanced (2021 Series) - Test 06 (22 Apr 2021)-img/p1_phy_s2_para7_8_1.png}{120.5}{0}
\qfigdef{p1_phy_s2_para9_10_1}{JEE Advanced (2021 Series) - Test 06 (22 Apr 2021)-img/p1_phy_s2_para9_10_1.png}{183.3}{0}
\qfigdef{p1_phy_s2_q09_A}{JEE Advanced (2021 Series) - Test 06 (22 Apr 2021)-img/p1_phy_s2_q09_A.png}{69.4}{0}
\qfigdef{p1_phy_s2_q09_B}{JEE Advanced (2021 Series) - Test 06 (22 Apr 2021)-img/p1_phy_s2_q09_B.png}{69.4}{0}
\qfigdef{p1_phy_s2_q09_C}{JEE Advanced (2021 Series) - Test 06 (22 Apr 2021)-img/p1_phy_s2_q09_C.png}{69.4}{0}
\qfigdef{p1_phy_s2_q09_D}{JEE Advanced (2021 Series) - Test 06 (22 Apr 2021)-img/p1_phy_s2_q09_D.png}{69.4}{0}
\qfigdef{p1_phy_s3_q02_1}{JEE Advanced (2021 Series) - Test 06 (22 Apr 2021)-img/p1_phy_s3_q02_1.png}{106.0}{0}
\qfigdef{p1_phy_s3_q03_1}{JEE Advanced (2021 Series) - Test 06 (22 Apr 2021)-img/p1_phy_s3_q03_1.png}{135.7}{0}
\qfigdef{p1_phy_s3_q06_1}{JEE Advanced (2021 Series) - Test 06 (22 Apr 2021)-img/p1_phy_s3_q06_1.png}{103.6}{0}
\qfigdef{p1_phy_s3_q07_1}{JEE Advanced (2021 Series) - Test 06 (22 Apr 2021)-img/p1_phy_s3_q07_1.png}{154.0}{0}
\qfigdef{p1_chem_s2_para7_8_1}{JEE Advanced (2021 Series) - Test 06 (22 Apr 2021)-img/p1_chem_s2_para7_8_1.png}{139.6}{0}
\qfigdef{p1_chem_s2_para9_10_1}{JEE Advanced (2021 Series) - Test 06 (22 Apr 2021)-img/p1_chem_s2_para9_10_1.png}{437.3}{1}
\qfigdef{p1_chem_s2_q09_A}{JEE Advanced (2021 Series) - Test 06 (22 Apr 2021)-img/p1_chem_s2_q09_A.png}{77.6}{0}
\qfigdef{p1_chem_s2_q09_B}{JEE Advanced (2021 Series) - Test 06 (22 Apr 2021)-img/p1_chem_s2_q09_B.png}{55.8}{0}
\qfigdef{p1_chem_s2_q09_C}{JEE Advanced (2021 Series) - Test 06 (22 Apr 2021)-img/p1_chem_s2_q09_C.png}{79.8}{0}
\qfigdef{p1_chem_s2_q09_D}{JEE Advanced (2021 Series) - Test 06 (22 Apr 2021)-img/p1_chem_s2_q09_D.png}{91.8}{0}
\qfigdef{p1_chem_s2_q10_A}{JEE Advanced (2021 Series) - Test 06 (22 Apr 2021)-img/p1_chem_s2_q10_A.png}{78.1}{0}
\qfigdef{p1_chem_s2_q10_B}{JEE Advanced (2021 Series) - Test 06 (22 Apr 2021)-img/p1_chem_s2_q10_B.png}{76.4}{0}
\qfigdef{p1_chem_s2_q10_4}{JEE Advanced (2021 Series) - Test 06 (22 Apr 2021)-img/p1_chem_s2_q10_4.png}{76.4}{0}
\qfigdef{p1_chem_s2_q10_D}{JEE Advanced (2021 Series) - Test 06 (22 Apr 2021)-img/p1_chem_s2_q10_D.png}{78.7}{0}
\qfigdef{p1_chem_s3_q01_1}{JEE Advanced (2021 Series) - Test 06 (22 Apr 2021)-img/p1_chem_s3_q01_1.png}{55.5}{0}
\qfigdef{p2_phy_s1_q01_1}{JEE Advanced (2021 Series) - Test 06 (22 Apr 2021)-img/p2_phy_s1_q01_1.png}{115.8}{0}
\qfigdef{p2_phy_s1_q02_1}{JEE Advanced (2021 Series) - Test 06 (22 Apr 2021)-img/p2_phy_s1_q02_1.png}{166.9}{0}
\qfigdef{p2_phy_s1_q05_1}{JEE Advanced (2021 Series) - Test 06 (22 Apr 2021)-img/p2_phy_s1_q05_1.png}{141.6}{0}
\qfigdef{p2_phy_s1_q06_1}{JEE Advanced (2021 Series) - Test 06 (22 Apr 2021)-img/p2_phy_s1_q06_1.png}{245.2}{0}
\qfigdef{p2_phy_s2_q08_1}{JEE Advanced (2021 Series) - Test 06 (22 Apr 2021)-img/p2_phy_s2_q08_1.png}{112.5}{0}
\qfigdef{p2_phy_s2_q09_1}{JEE Advanced (2021 Series) - Test 06 (22 Apr 2021)-img/p2_phy_s2_q09_1.png}{146.6}{0}
\qfigdef{p2_phy_s2_q10_1}{JEE Advanced (2021 Series) - Test 06 (22 Apr 2021)-img/p2_phy_s2_q10_1.png}{64.8}{0}
\qfigdef{p2_phy_s2_q10_2}{JEE Advanced (2021 Series) - Test 06 (22 Apr 2021)-img/p2_phy_s2_q10_2.png}{101.9}{0}
\qfigdef{p2_phy_s2_q10_3}{JEE Advanced (2021 Series) - Test 06 (22 Apr 2021)-img/p2_phy_s2_q10_3.png}{84.2}{0}
\qfigdef{p2_phy_s2_q10_4}{JEE Advanced (2021 Series) - Test 06 (22 Apr 2021)-img/p2_phy_s2_q10_4.png}{146.8}{0}
\qfigdef{p2_phy_s3_q01_1}{JEE Advanced (2021 Series) - Test 06 (22 Apr 2021)-img/p2_phy_s3_q01_1.png}{180.6}{0}
\qfigdef{p2_phy_s3_q04_1}{JEE Advanced (2021 Series) - Test 06 (22 Apr 2021)-img/p2_phy_s3_q04_1.png}{118.3}{0}
\qfigdef{p2_phy_s3_q07_1}{JEE Advanced (2021 Series) - Test 06 (22 Apr 2021)-img/p2_phy_s3_q07_1.png}{183.5}{0}
\qfigdef{p2_phy_s3_q08_1}{JEE Advanced (2021 Series) - Test 06 (22 Apr 2021)-img/p2_phy_s3_q08_1.png}{92.9}{0}
\qfigdef{p2_chem_s1_q01_1}{JEE Advanced (2021 Series) - Test 06 (22 Apr 2021)-img/p2_chem_s1_q01_1.png}{178.9}{0}
\qfigdef{p2_chem_s1_q01_2}{JEE Advanced (2021 Series) - Test 06 (22 Apr 2021)-img/p2_chem_s1_q01_2.png}{91.2}{0}
\qfigdef{p2_chem_s1_q01_3}{JEE Advanced (2021 Series) - Test 06 (22 Apr 2021)-img/p2_chem_s1_q01_3.png}{58.0}{0}
\qfigdef{p2_chem_s1_q02_1}{JEE Advanced (2021 Series) - Test 06 (22 Apr 2021)-img/p2_chem_s1_q02_1.png}{81.2}{0}
\qfigdef{p2_chem_s1_q02_2}{JEE Advanced (2021 Series) - Test 06 (22 Apr 2021)-img/p2_chem_s1_q02_2.png}{72.5}{0}
\qfigdef{p2_chem_s1_q02_3}{JEE Advanced (2021 Series) - Test 06 (22 Apr 2021)-img/p2_chem_s1_q02_3.png}{69.6}{0}
\qfigdef{p2_chem_s1_q02_4}{JEE Advanced (2021 Series) - Test 06 (22 Apr 2021)-img/p2_chem_s1_q02_4.png}{65.6}{0}
\qfigdef{p2_chem_s1_q03_1}{JEE Advanced (2021 Series) - Test 06 (22 Apr 2021)-img/p2_chem_s1_q03_1.png}{73.3}{0}
\qfigdef{p2_chem_s1_q03_A}{JEE Advanced (2021 Series) - Test 06 (22 Apr 2021)-img/p2_chem_s1_q03_A.png}{33.2}{0}
\qfigdef{p2_chem_s1_q03_B}{JEE Advanced (2021 Series) - Test 06 (22 Apr 2021)-img/p2_chem_s1_q03_B.png}{91.6}{0}
\qfigdef{p2_chem_s1_q03_C}{JEE Advanced (2021 Series) - Test 06 (22 Apr 2021)-img/p2_chem_s1_q03_C.png}{41.6}{0}
\qfigdef{p2_chem_s1_q03_D}{JEE Advanced (2021 Series) - Test 06 (22 Apr 2021)-img/p2_chem_s1_q03_D.png}{44.9}{0}
\qfigdef{p2_chem_s1_q05_D}{JEE Advanced (2021 Series) - Test 06 (22 Apr 2021)-img/p2_chem_s1_q05_D.png}{73.1}{0}
\qfigdef{p2_chem_s2_q08_1}{JEE Advanced (2021 Series) - Test 06 (22 Apr 2021)-img/p2_chem_s2_q08_1.png}{53.0}{0}
\qfigdef{p2_chem_s2_q08_2}{JEE Advanced (2021 Series) - Test 06 (22 Apr 2021)-img/p2_chem_s2_q08_2.png}{33.2}{0}
\qfigdef{p2_chem_s2_q08_3}{JEE Advanced (2021 Series) - Test 06 (22 Apr 2021)-img/p2_chem_s2_q08_3.png}{34.3}{0}
\qfigdef{p2_chem_s3_q05_1}{JEE Advanced (2021 Series) - Test 06 (22 Apr 2021)-img/p2_chem_s3_q05_1.png}{111.5}{0}

\begin{document}
\begin{center}{\Large\bfseries JEE Advanced (2021 Series) --- Test 06 (22 Apr 2021)}\end{center}
\begin{multicols}{2}
\papersec{1}
\partsec{1}{PHYSICS}
\sect{Section-I(i)}
\qn{1} A small solid ball of density $\rho_0$ is held inside at point $A$ in a closed cubical container of side $L$, filled with an ideal liquid of density $4\rho_0$ as shown in the figure. Now, if the container starts moving with constant acceleration $a_0$ horizontally and the ball is released from point $A$ simultaneously, It has been observed that the ball hits the point $Q$ after time of $T_0$. Points $A$, $P$, $Q$, $R$ and $S$ are lying in the same plane. Then choose the \textbf{CORRECT} option(s) :
\qfig{p1_phy_s1_q01_1}
\begin{qoptions}
\opt{A}{For ball to hit the top of container at end $Q$, $a_0=\dfrac{3g}{2}$}
\opt{B}{For ball to hit the top of container at end $Q$, $a_0=2g$}
\opt{C}{Ball hits the top of container at end $Q$ after a time $T_0=\sqrt{\dfrac{L}{3g}}$}
\opt{D}{Ball hits the top of container at end $Q$ after a time $T_0=\sqrt{\dfrac{2L}{3g}}$}
\end{qoptions}
\qn{2} In a radioactive decay process $A$ decays to $B$, $A\to B$. Two graphs of number of nuclei of $A$ and $B$ versus time is given. Then choose the \textbf{CORRECT} option(s)
\qfig{p1_phy_s1_q02_1}
\qfig{p1_phy_s1_q02_2}
\begin{qoptions}
\opt{A}{$t_2-t_1=4$}
\opt{B}{$t_2-t_1=2$}
\opt{C}{$t_1=2\log_25$}
\opt{D}{$t_2=\log_2100$}
\end{qoptions}
\qn{3} A non-uniform rope of length $L$ is placed on the ground and Rope is pulled by constant force $F_0$ as shown in the diagram, there is no friction anywhere, linear mass density of rope is given by $\lambda_x=\lambda_0\left(1+\dfrac xL\right)$, where $x$ is measured from the end $B$. If tension at a distance of $x$ from $B$ is $T_x$, total mass of the rope is $M$ and acceleration of rope is $a_0$. Then choose the \textbf{CORRECT} option(s).
\qfig{p1_phy_s1_q03_1}
\begin{qoptions}
\opt{A}{$\displaystyle\int_0^LT_x\cdot dx=\dfrac{4F_0L}{9}$}
\opt{B}{$M=\dfrac{3\lambda_0L}{2}$}
\opt{C}{$a_0=\dfrac{2F_0}{3\lambda_0L}$}
\opt{D}{$T_x=\dfrac{2F_0}{3L}\left[x+\dfrac{x^2}{2L}\right]$}
\end{qoptions}
\qn{4} A non conducting uniform thin shell of mass $m$ and radius $R$ is having uniform charge density $+\sigma$ on one half and $-\sigma$ on another half as shown in figure. It is placed on a rough non conducting horizontal plane. At $t=0$ a uniform electric field $\vec E=-E_0\hat j$ N/C is switched on and the solid sphere starts rolling without slipping. Then which of the following is/are \textbf{CORRECT}.
\qfig{p1_phy_s1_q04_1}
\begin{qoptions}
\opt{A}{The acceleration of the topmost point of the hollow sphere at $t=0$, is $\dfrac{36}{25}\dfrac{\pi\sigma E_0R^2}{m}$.}
\opt{B}{The magnitude of frictional force at $t=0$ is $\dfrac{18}{25}\pi\sigma E_0R^2$.}
\opt{C}{The acceleration of centre of mass of the hollow sphere at $t=0$, is $\dfrac{12}{25}\dfrac{\pi\sigma E_0R^2}{m}$.}
\opt{D}{The magnitude of frictional force at $t=0$ is $\dfrac{9}{25}\pi\sigma E_0R^2$.}
\end{qoptions}
\qn{5} A point object is placed at $30$ cm from a convex glass lens $\left(\mu_g=\dfrac32\right)$ of focal length $20$ cm. The final image of object will be formed at infinity if -
\begin{qoptions}
\opt{A}{Another concave lens of focal length $60$ cm is placed in contact with the previous lens}
\opt{B}{Another convex lens of focal length $60$ cm is placed at a distance of $30$ cm from the first lens}
\opt{C}{The whole system is immersed in a liquid of refractive index $4/3$}
\opt{D}{The whole system is immersed in a liquid of refractive index $9/8$}
\end{qoptions}
\qn{6} In the cylindrical region of radius $R=2$ m, there exists a time varying magnetic field $B$ such that $\dfrac{dB}{dt}=2$ Tesla/sec. A charged particle having charge $q=2$ C is placed at the point $P$ at a distance $d$ from its center $O$. Now, the particle is moved in the direction perpendicular to the line $OP$ by an external agent up to infinity so that there is no gain in kinetic energy of charged particle. If ($d>R$), then which of the folllowing options is/are \textbf{CORRECT}?
\qfig{p1_phy_s1_q06_1}
\begin{qoptions}
\opt{A}{work done by external agent is $4\pi$ Joule when $d=4$ m}
\opt{B}{work done by external agent is $4\pi$ Joule when $d=8$ m.}
\opt{C}{work done by external agent is independent of $d$.}
\opt{D}{work done by external agent is positive.}
\end{qoptions}
\sect{Section-I(ii)}
\paraq{7}{8}
A uniform wire frame of linear mass density $\lambda$ having three sides each of length $2a$ is kept on a smooth horizontal surface. An impulse $J$ is applied at one end as shown in the figure. $P$ is the midpoint of $AB$.
\qfig{p1_phy_s2_para7_8_1}
\qn{7} The angular velocity of system just after the impulse is
\begin{qoptions}
\opt{A}{$\dfrac{3J}{22\lambda a^2}$}
\opt{B}{$\dfrac{3J}{11\lambda a^2}$}
\opt{C}{$\dfrac{J}{11\lambda a^2}$}
\opt{D}{$\dfrac{2J}{11\lambda a^2}$}
\end{qoptions}
\qn{8} The velocity of point $P$ just after the impulse is
\begin{qoptions}
\opt{A}{$\dfrac{J}{\lambda a}\hat j$}
\opt{B}{$\dfrac{J}{6\lambda a}\hat j$}
\opt{C}{$\dfrac{J}{\lambda a}\left(\dfrac{2}{11}\hat i+\dfrac16\hat j\right)$}
\opt{D}{$\dfrac{J}{\lambda a}\left(\dfrac{1}{11}\hat i+\dfrac16\hat j\right)$}
\end{qoptions}
\paraq{9}{10}
Consider YDSE arrangement shown. The screen in the arrangement starts accelerating from rest in the positive $x$-direction with acceleration $a=Kt$ at time $t=0$ (here $K$ is a constant and $t$ is time)
\qfig{p1_phy_s2_para9_10_1}
\qn{9} Rate of change of fringe width ($R$) for the fringe pattern on the screen varies with time as
\begin{qoptions}
\opt{A}{\qfigi{p1_phy_s2_q09_A}}
\opt{B}{\qfigi{p1_phy_s2_q09_B}}
\opt{C}{\qfigi{p1_phy_s2_q09_C}}
\opt{D}{\qfigi{p1_phy_s2_q09_D}}
\end{qoptions}
\qn{10} Velocity of $n^{\text{th}}$ maxima (above ``$O$'') at any time $t$ in ground frame is -
\begin{qoptions}
\opt{A}{$\dfrac{Kt^2}{2}\left(\hat i+\dfrac{n\lambda}{d}\hat j\right)$}
\opt{B}{$Kt\left(\hat i+\dfrac{n\lambda}{d}\hat j\right)$}
\opt{C}{$\dfrac{Kt^2}{2}\left(\dfrac{n\lambda}{d}\hat i+\hat j\right)$}
\opt{D}{None of these}
\end{qoptions}
\sect{Section-II}
\qn{1} In a radioacitve decay $A$ is converting into $B$ and $B$ is converting into $C$. It is given that at $t=0$ nuclei of $B$ and $C$ are not present, decay constant for $A$ to $B$ is $\lambda_1$ and decay constant for $B$ to $C$ is $\lambda_2$, take $\lambda_1=5\lambda$ and $\lambda_2=2\lambda$. The time after which $B$ will have maximum No. of nuclei is of the form $\dfrac{\ln(5/2)}{2B_0\lambda}$. The value of $B_0$ will be :-
\[A\xrightarrow{\ \lambda_1\ }B\xrightarrow{\ \lambda_2\ }C\]
\qn{2} A system of two planks and a sphere of radius $R$ is in motion as shown in figure. Radius of the sphere is $R$ and there is no slipping anywhere. It is given that $R\alpha=2a_0$ where $\alpha$ is angular acceleration of sphere and acceleration of upper block is $a_1=ka_0$ where $k$ is a $+$ve constant then the value of $k$ will be :
\qfig{p1_phy_s3_q02_1}
\qn{3} An infinite uniform current carrying wire is kept along $z$-axis, carrying current $I_0$ in the direction of the positive $z$-axis. $OABCDEFG$ represents a circle (where all the points are equally spaced) whose centre at point $(4\text{m},0\text{m})$ and radius $4$ m as shown in the figure. $\displaystyle\int_{DEF}\vec B\cdot\vec{d\ell}=\dfrac{\mu_0I_0k}{2}$ in S.I. unit, then the value of $k$ is :
\qfig{p1_phy_s3_q03_1}
\qn{4} A closed organ pipe of length $L$ and an open organ pipe contain gases of densities $\rho_1$ and $\rho_2$ respectively. The compressibility of gases are equal in both the pipes. Both the pipes are vibrating in their first overtone with same frequency. The length of open organ pipe is $kL\left(\dfrac{\rho_1}{\rho_2}\right)^{1/2}$. Find value of $k$
\qn{5} A non-conducting thin disc of radius $R$ charged uniformly over one side rotates uniformly with an angular velocity $\omega$ about an axis passing through its periphery point $P$ and perpendicular to its plane. The electrostatic potential at the point $P$ due to charges on the disc is $V_0$ before the rotation. The magnetic field at $P$ due to the moving charges upon rotation is $2k\mu_0\epsilon_0\omega V_0$. Find the value of $k$
\qn{6} A uniform solid cylinder of mass `$m$' rests on two horizontal planks. A thread is wound on the cylinder. The hanging end of the thread is pulled vertically down with a constant force $F$, which still does not bring any sliding of the cylinder. If the coefficient of friction between the cylinder and planks in equal to $\mu$. What is the maximum acceleration $a_{\max}$ of the axis of cylinder rolling over the planks. (Assume $\mu=1/3$ and take $g=10$ m/sec.$^2$)
\qfig{p1_phy_s3_q06_1}
\qn{7} A capacitor of capacity $6\ \mu$F and initial charge $160\ \mu$C is connected with a key $S$ and resistance as shown in figure. Point $M$ is earthed. If key is closed at $t=0$ then the current through resistance $R=1\ \Omega$ at $t=16\ \mu$ sec will be \underline{\hspace{1cm}} Amp. (Assume $e=2.7$)
\qfig{p1_phy_s3_q07_1}
\qn{8} The distance between the parallel plates of a isolated charged condenser is $d=5$ cm and the intensity of the field $E=300$ V cm$^{-1}$. A slab of dielectric constant $K=5$ and $1$ cm wide is inserted parallel to the plates. Determine the potential difference (in volt) between the plates, after the slab is inserted.
\partsec{2}{CHEMISTRY}
\sect{Section-I(i)}
\qn{1} Select the incorrect statement(s) among given
\begin{qoptions}
\opt{A}{The highest occupied molecular orbital (HOMO) in $\mathrm{N_2}$ has no s-character.}
\opt{B}{Anhydrous $\mathrm{CaCl_2}$ can be used for drying of $\mathrm{NH_3}$.}
\opt{C}{Pt can form more stable complex with $+4$ oxidation state as compare to $\mathrm{Ni^{+4}}$.}
\opt{D}{$\mathrm{(O_2)[PF_6]}$ is paramagnetic.}
\end{qoptions}
\qn{2} Find correct statements for Coordiantion compounds.
\begin{qoptions}
\opt{A}{$\mathrm{[PtCl_2(NH_3)_2]}$ \& $\mathrm{[Cu(NH_3)_4]^{2+}}$ are square planar complexes.}
\opt{B}{Complex $\mathrm{[Fe(C_2O_4)_3]^{3-}}$, does not exhibit geometrical isomerism while exhibits optical isomerism.}
\opt{C}{CO bond order : $\mathrm{[Ni(CO)_4]}$ $>$ $\mathrm{[Co(CO)_4]^-}$ $>$ $\mathrm{[Fe(CO)_4]^{2-}}$.}
\opt{D}{$\mathrm{[Ni(Cl_2)(PPh_3)_2]}$ has two geometrical isomers.}
\end{qoptions}
\qn{3} Equivalent weight of the underlined ($\mathrm{HNO_3}$ as oxidising agent and HCl as reducing agent) substances in the given reaction is respectively (approx).\\
(i) $\mathrm{Cu+\underline{HNO_3}\longrightarrow Cu(NO_3)_2+NO+H_2O}$\\
(ii) $\mathrm{K_2Cr_2O_7+\underline{HCl}\longrightarrow KCl+CrCl_3+Cl_2+H_2O}$
\begin{qoptions}
\opt{A}{$63,\ 36.5$}
\opt{B}{$63,\ 73$}
\opt{C}{$84,\ 85$}
\opt{D}{$84,\ 42.5$}
\end{qoptions}
\qn{4} The radial distribution functions $[P(r)]$ is used to determine the most probable radius, which is used to find the electron in a given orbital $\dfrac{dP(r)}{dr}$ for 1s--orbital of hydrogen like atom having atomic number $Z$, is
\[\dfrac{dP(r)}{dr}=\dfrac{16\pi Z^3}{a_0^3}\left(2r-\dfrac{2Zr^2}{a_0}\right)e^{-2Zr/a_0}\]
Which of the following statement(s) is/are correct ?
\begin{qoptions}
\opt{A}{At the point of maximum value of radial distribution function, $\dfrac{dP(r)}{dr}=0$; one antinode is present}
\opt{B}{Most probable radius of $\mathrm{Li^{2+}}$ is $\dfrac{a_0}{3}$}
\opt{C}{Most probable radius of $\mathrm{He^{+}}$ is $\dfrac{a_0}{2}$}
\opt{D}{Most probable radius of hydrogen atom is $a_0$}
\end{qoptions}
\qn{5} Consider the following table and selct correct option/s.
\begin{center}
{\footnotesize\setlength\tabcolsep{3pt}\begin{tabular}{@{}p{.2\linewidth}p{.3\linewidth}p{.17\linewidth}p{.22\linewidth}@{}}
\toprule
Compound & bond dissociation energy (kJ/mole) & Bond angle & Boiling point ($^\circ$C)\\
\midrule
$\mathrm{H_2O}$ & $a_1$ & $b_1$ & $c_1$\\
$\mathrm{H_2S}$ & $a_2$ & $b_2$ & $c_2$\\
$\mathrm{H_2Se}$ & $a_3$ & $b_3$ & $c_3$\\
$\mathrm{H_2Te}$ & $a_4$ & $b_4$ & $c_4$\\
\bottomrule
\end{tabular}}
\end{center}
\begin{qoptions}
\opt{A}{$a_1>a_2>a_3>a_4$}
\opt{B}{$b_1>b_2>b_3>b_4$}
\opt{C}{$c_1>c_2>c_3>c_4$}
\opt{D}{$c_2<c_3<c_4<c_1$}
\end{qoptions}
\qn{6} How many of the following statement is/are \textbf{INCORRECT}?
\begin{qoptions}
\opt{A}{EAN of $\mathrm{[Rh(ox)_3]^{3-}}$ is $56$}
\opt{B}{$d_{\mathrm{C-C}}$ in Zeise's salt is less than that of $\mathrm{C_2F_4}$ molecule}
\opt{C}{Red bauxite consists of $\mathrm{Fe_2O_3}$ as minor impurity}
\opt{D}{In Hall-Heroult process, the produced Al floats at the top of the electrolyte}
\end{qoptions}
\sect{Section-I(ii)}
\paraq{7}{8}
The cell potential for the chemical reaction :
\begin{gather*}\mathrm{Hg_2^{2+}(aq)+NO_3^-(aq)+3H_3O^+(aq)}\\ \mathrm{\to2Hg^{2+}(aq)+HNO_2(aq)+4H_2O(\ell)}\end{gather*}
is measured under standard conditions in the electrochemical cell shown in the accompanying diagram.
\qfig{p1_chem_s2_para7_8_1}
\qn{7} How many moles of electron pass through the circuit when $0.60$ mole of $\mathrm{Hg^{2+}}$ and $0.30$ mole of $\mathrm{HNO_2}$ are produced in the cell that contain $0.50$ mole of $\mathrm{Hg_2^{\,2+}}$ and $0.40$ mole of $\mathrm{NO_3^-}$ at the beginning of the reaction ?
\begin{qoptions}
\opt{A}{$0.30$}
\opt{B}{$0.60$}
\opt{C}{$0.15$}
\opt{D}{$1.20$}
\end{qoptions}
\qn{8} How long will it take to produce $0.10$ mole of $\mathrm{HNO_2}$ by this reaction if the cell delivers a constant current of $10$ A ?
\begin{qoptions}
\opt{A}{$965$ s}
\opt{B}{$193$ s}
\opt{C}{$1930$ s}
\opt{D}{$482.5$ s}
\end{qoptions}
\qfullbegin
\paraq{9}{10}
\qfig{p1_chem_s2_para9_10_1}
\qfullend
\qn{9} The compound (D) is
\begin{qoptions}
\opt{A}{\qfigi{p1_chem_s2_q09_A}}
\opt{B}{\qfigi{p1_chem_s2_q09_B}}
\opt{C}{\qfigi{p1_chem_s2_q09_C}}
\opt{D}{\qfigi{p1_chem_s2_q09_D}}
\end{qoptions}
\qn{10} The compound F is
\begin{qoptions}
\opt{A}{\qfigi{p1_chem_s2_q10_A}}
\opt{B}{\qfigi{p1_chem_s2_q10_B}}
\opt{C}{\qfigi{p1_chem_s2_q10_4}}
\opt{D}{\qfigi{p1_chem_s2_q10_D}}
\end{qoptions}
\sect{Section-II}
\qn{1} \qfigi{p1_chem_s3_q01_1}$\ \xrightarrow{\text{NBS}}$ Number of monobromoproducts formed (including stereoisomerism)
\qn{2} Following data is given for two ideal gas samples taken in two seperate containers :
\begin{center}
{\footnotesize\setlength\tabcolsep{3pt}\begin{tabular}{@{}p{.1\linewidth}p{.17\linewidth}p{.15\linewidth}p{.2\linewidth}p{.25\linewidth}@{}}
\toprule
Gas & Pressure & Volume & Temperature & Mean free paths($\lambda$)\\
\midrule
X & $2$ atm & $V$ litre & $800$ K & $0.08$ nm\\
Y & $8$ atm & $V$ litre & $200$ K & $0.02$ nm\\
\bottomrule
\end{tabular}}
\end{center}
\& $\dfrac{(Z_{11})_X}{(Z_{11})_Y}=\dfrac1{16}$\\
Find the ratio of rate of diffusion of gas $X$ and gas $Y$ in similar condition of pressure, temperature, etc.
\qn{3} When $\mathrm{P_4}$ reacts with NaOH \& gives a gas (A) \& forms another compound (B) containing phosphorous. Find out the difference in oxidation state of Phosphorous in (A) and (B). (Consider A and B as a major product)
\qn{4} (i) $\mathrm{Ti^{+3}}\Rightarrow$ Purple \quad (ii) $\mathrm{Fe^{+2}}\Rightarrow$ Green\\
(iii) $\mathrm{Fe^{+3}}\Rightarrow$ Yellow \quad (iv) $\mathrm{Ni^{+2}}\Rightarrow$ Green\\
(v) $\mathrm{Cu^{+2}}\Rightarrow$ Blue \quad (vi) $\mathrm{Sc^{3+}}\Rightarrow$ Colourless\\
(vii) $\mathrm{Ti^{4+}}\Rightarrow$ Colourless \quad (viii) $\mathrm{Zn^{2+}}\Rightarrow$ Colourless\\
How many colour of aqueous cation are correctly matched.
\qn{5} In the saturated aqueous solution of $\mathrm{PbCl_2}$, the freezing point of water decreases by $\left(\dfrac{D}{100}\right)^\circ$C, then `D' is :\\
(Given $k_{sp}$ of $\mathrm{PbCl_2}=4\times10^{-6}$, $K_{f,\,\mathrm{water}}=2$ K kg/mole, assume molality $=$ molarity)
\qn{6} A reversible cyclic process involves $6$ steps. In step -1 and 3, system absorb $500$ J and $800$ J of heat from a heat reservoir at temperature $250$K \& $200$K, respectively. Step 2, 4, 6, are adiabatic such that the temperature of one reservoir changes to that of next. Total work done by the system in whole cycle is $700$ J. Find the temperature (in K) during step 5 if it exchange heat from a reservoir at temperature $T_5$
\qn{7} Various isomeric alkynes have molecular formula $\mathrm{C_3FClBrI}$. For such alkynes if\\
$p=$ Number of positional isomers\\
$q=$ Number of optically active isomers\\
$r=$ Number of fractions obtained on fractional distillation of a mixture of all possible alkynes\\
$s=$ Number of diastereomeric pairs\\
Then find the value of $\left(\dfrac{p+q-r+s}{5}\right)$.
\qn{8} An optically active substance `A' decomposes into optically active substances `B' and `C' as follows
\[\mathrm{A}\xrightarrow{K=0.001\,\mathrm{min}^{-1}}\mathrm{2B+C}\]
The specific rotations of A, B and C are $+40^\circ$, $+10^\circ$ and $-30^\circ$ per mole, respectively. If initially A and C were present in $4:3$ mole ratio, the time (in min), after which the sample becomes optically inactive is ($\ln2=0.7$, $\ln5=1.6$, $\ln7=2.0$, $\ln13=2.5$)
\partsec{3}{MATHEMATICS}
\sect{Section-I(i)}
\qn{1} Consider two sets
\begin{gather*}A=\left\{\dfrac pq;\ p,q\text{ are coprime odd integers}\right\}\\ \text{and}\quad B=\left\{\dfrac{p+1}{q}:p,q\text{ odd integers},\ p\neq-1\right\}.\end{gather*}
Functions $f:A\to\mathbb R$ \& $g:B\to\mathbb R$ are both defined as $f(x)=x^x$ \& $g(x)=x^{|x|}$ on respective domains. Then :
\begin{qoptions}
\opt{A}{$f\!\left(\dfrac13\right)\cdot g\!\left(\dfrac{-2}{3}\right)=\dfrac{\sqrt[3]{4}}{3}$}
\opt{B}{$g(x)$ is an even function.}
\opt{C}{$f(1)+g(2)=5$}
\opt{D}{$\left[g\!\left(\dfrac{-6}{5}\right)\right]^{5/2}-\left[g\!\left(\dfrac45\right)\right]^{15/4}-\left[f\!\left(\dfrac35\right)\right]^5=1$}
\end{qoptions}
\qn{2} If sides of triangle $ABC$ are $a$, $b$ and $c$ such that $2b=a+c$ then
\begin{qoptions}
\opt{A}{$\dfrac bc>\dfrac23$}
\opt{B}{$\dfrac bc>2$}
\opt{C}{$\dfrac bc<2$}
\opt{D}{$\dfrac bc<\dfrac32$}
\end{qoptions}
\qn{3} If $z_1,z_2,z_3$ are any three roots of the equation $z^6=(z+1)^6$, then $\arg\left(\dfrac{z_1-z_2}{z_2-z_3}\right)$ can be equal to
\begin{qoptions}
\opt{A}{$0$}
\opt{B}{$\pi$}
\opt{C}{$\dfrac\pi4$}
\opt{D}{$-\dfrac\pi4$}
\end{qoptions}
\qn{4} If foci of an ellipse be $(-1,2)$ and $(-2,3)$ and its tangent at a point $A$ is $2x+3y+9=0$
\begin{qoptions}
\opt{A}{Length of the minor axis of the ellipse will be $2\sqrt{14}$}
\opt{B}{Co-ordinate of the point `$A$' will be $\left(\dfrac{-32}{9},\dfrac{-17}{27}\right)$}
\opt{C}{Distance between the foci is $2\sqrt2$}
\opt{D}{Product of the perpendiculars from foci to any tangent is $56$}
\end{qoptions}
\qn{5} Let $A_1,A_2,A_3,A_4$ are in G.P. If the H.M. of $A_1,A_2$ is $12$ and that of $A_2,A_3$ is $36$, then
\begin{qoptions}
\opt{A}{$3A_1-A_2=4$}
\opt{B}{$9A_1-A_3=0$}
\opt{C}{$A_4-4A_2-2A_3=-24$}
\opt{D}{The sum of digits in $A_4$ in $6$}
\end{qoptions}
\qn{6} Let the function \[f(x)=\begin{cases}2x-x^3+\ln(a^2-3a+3), & 0\le x<2\\ x-6, & x\ge2\end{cases}\] It possess a local minimum at $x=2$ for
\begin{qoptions}
\opt{A}{$a=0$}
\opt{B}{$a=3$}
\opt{C}{$a=4$}
\opt{D}{$a=-1$}
\end{qoptions}
\sect{Section-I(ii)}
\paraq{7}{8}
Consider the quadratic equation $2x^2-(4m+2)x+m^2+m=0$, $m\in\mathbb R$
\qn{7} The number of positive integer values of `$m$' such that the equation has exactly one root in $(2,3)$ is
\begin{qoptions}
\opt{A}{$3$}
\opt{B}{$4$}
\opt{C}{$5$}
\opt{D}{$6$}
\end{qoptions}
\qn{8} The number of negative integral values of `$m$' such that $m>-10$ and atleast one root of the equation is smaller than `$2$' is
\begin{qoptions}
\opt{A}{$8$}
\opt{B}{$9$}
\opt{C}{$6$}
\opt{D}{$4$}
\end{qoptions}
\paraq{9}{10}
Consider a standard ellipse whose focus and corresponding foot of directrix are $\left(\sqrt7,0\right)$ and $\left(\dfrac{16}{\sqrt7},0\right)$ and a circle with equation $x^2+y^2=r^2$. If in the first quadrant, the common tangent to a circle of this family and the above ellipse meets the co-ordinate axes in $A$ and $B$ then
\qn{9} If `$P$' is a variable point on the above ellipse with $S$, $S'$ are foci then the maximum possible area of $\triangle PSS'$ is
\begin{qoptions}
\opt{A}{$\sqrt7$ sq. units}
\opt{B}{$2\sqrt7$ sq. units}
\opt{C}{$3\sqrt7$ sq. units}
\opt{D}{$4\sqrt7$ sq. units}
\end{qoptions}
\qn{10} Slope of line joining $AB$ when $r=\dfrac72$ is
\begin{qoptions}
\opt{A}{$-\sqrt{\dfrac{13}{15}}$}
\opt{B}{$\sqrt{\dfrac{13}{15}}$}
\opt{C}{$\sqrt{\dfrac{15}{13}}$}
\opt{D}{$-\sqrt{\dfrac{15}{13}}$}
\end{qoptions}
\sect{Section-II}
\qn{1} Given : $S_n=\displaystyle\sum_{r=1}^nt_r=\dfrac{n(2n^2+9n+13)}{6}$ Find $\displaystyle\sum_{r=1}^\infty\dfrac{1}{r\cdot\sqrt{t_r}}$
\qn{2} Let $I(n)=\displaystyle\int_1^ex^3\cdot(\ln x)^n\,dx$, where $n\in\mathbb N$. Find the value of $\ln(4I(5)+5I(4))$.
\qn{3} Consider a sequence $f_n(x)$ defined as $f_1(x)=1$ \& $f_{n+1}(x)=\displaystyle\int_0^\infty\dfrac{e^{-tx}}{f_n(t)}\,dt$. Evaluate the value of $\dfrac{f_8(2)}{f_6(2)}$.
\qn{4} Let $A=\begin{bmatrix}3&1\\0&2\end{bmatrix}$. If $A^3-5A^2+9A=B$. Find trace of matrix $B$.
\qn{5} If $f(x)=\dfrac{ax+b}{bx-a}$, $x\neq\dfrac ab$, $a>0,b>0$ and if $g(x)=f(f(x))+f\!\left(f\!\left(\dfrac1x\right)\right)$ and $g:[1,\infty)\to[2,\infty)$ then the value of $x$ for which $g(x)=\dfrac52$, is
\qn{6} It is given that $\log_6a+\log_6b+\log_6c=6$, where $a$, $b$ and $c$ are positive integers that form an increasing geometric sequence and $b-a$ is the square of an integer. Find $a+b+c$.
\qn{7} Evaluate the value of the limit $L=\displaystyle\lim_{x\to0}\left(\dfrac{1}{\ln\left(x+\sqrt{x^2+1}\right)}-\dfrac{1}{\ln(x+1)}\right)$
\qn{8} Consider the curves $C_1:|z-2|=2+\mathrm{Re}(z)$ and $C_2:|z|=3$ (where $x,y\in\mathbb R$ and $i=\sqrt{-1}$). They intersect at $P$ and $Q$ is the first and fourth quadrants respectively. Tangents to $C_1$ at $P$ and $Q$ intersect the $x$-axis at $R$ and tangents to $C_2$ at $P$ and $Q$ intersect the $x$-axis at $S$. If area of $\triangle PRS$ is $\lambda\sqrt2$ sq.units, then the value of $(\lambda^2)$.
\papersec{2}
\partsec{1}{PHYSICS}
\sect{Section-I(i)}
\qn{1} A solid cone of base radius `$r$' and half angle $45^\circ$ made of material have density $\dfrac{3\rho}{2}$. Cone is placed inside a container which is filled with two immiscible liquid of density $\rho$ and $2\rho$ respectively. The cone is at rest with respect to container as shown in diagram. Which of the following is/are \textbf{CORRECT} statements ? (Neglect effect of atmospheric pressure)
\qfig{p2_phy_s1_q01_1}
\begin{qoptions}
\opt{A}{The ratio of volume of cone in liquid 1 and liquid 2 is 1.}
\opt{B}{The force exerted by the liquid 1 on the cone in vertical direction is $\dfrac{\pi r^2}{3}\rho g$}
\opt{C}{Option (A) and (B) both are correct.}
\opt{D}{Option (A) and (B) both are wrong.}
\end{qoptions}
\qn{2} Consider a square frame made of uniformly charged non conducting wires of length $d$ and linear charge density $\lambda\dfrac{\mathrm C}{\mathrm m}$. Choose the \textbf{CORRECT} statement(s).
\qfig{p2_phy_s1_q02_1}
\begin{qoptions}
\opt{A}{The electric field at $A$ is zero.}
\opt{B}{The potential at centre $A$ is higher than potential at $B$.}
\opt{C}{The net electric field at $B$ is $\left[\dfrac{32K\lambda}{\sqrt{150}d}\right]\dfrac{\mathrm N}{\mathrm C}$.}
\opt{D}{The potential at $B$ is $V=4K\lambda\log_e\left[\dfrac{\sec\theta+\tan\theta}{\sec\theta-\tan\theta}\right]$ $\dfrac{\mathrm J}{\mathrm C}$ : where $\theta$ is $\sin^{-1}\left[\dfrac{1}{\sqrt6}\right]$}
\end{qoptions}
\qn{3} In rural areas where electric power lines carry current to big cities, it is possible to generate small amount of electricity from the lines by means of induction in a single coil. The overhead power lines carry relatively large currents that periodically reverse in direction 50 time per second. How would you orient the plane of the coil to produce maximum induced current if the power lines ran in a north-to-south direction ?
\begin{qoptions}
\opt{A}{Parallel to the ground.}
\opt{B}{Perpendicular to the ground in the North-South directions.}
\opt{C}{Perpendicular to the ground in the East-West direction.}
\opt{D}{The orientation of the coil would not affect the size of the induced current.}
\end{qoptions}
\qn{4} The engines of a spacecraft in circular orbit about the earth are fired so as to give the craft a forward thrust. How will the space craft's energy be changed after the firing, assuming another circular orbit is achieved ?
\begin{qoptions}
\opt{A}{K.E. increased and P.E. decreased.}
\opt{B}{K.E. decreased and P.E. increased}
\opt{C}{Both K.E. and P.E. increased}
\opt{D}{Total energy increases.}
\end{qoptions}
\qn{5} A solid ball of mass $M$ and radius `$R$' is rolling without slipping on a flat horizontal surface at an initial angular velocity $\omega_0$. It hits a small bump of height $h$. Assume that the ball pivots about the tip of the bump during and after impact. Which of the following statements are \textbf{TRUE} ?
\qfig{p2_phy_s1_q05_1}
\begin{qoptions}
\opt{A}{The angular velocity of the ball immediately after impact is $\omega=\omega_0\left(1-\dfrac{5h}{7R}\right)$}
\opt{B}{The minimum initial angular velocity so that the ball just makes it over the bumb is $\omega_0=\dfrac{7}{(7R-5h)}\sqrt{\dfrac{10gh}{7}}$}
\opt{C}{The angular velocity of the ball immediately after impact is $\omega=\omega_0\left(1-\dfrac{3h}{7R}\right)$}
\opt{D}{During the impact gravitation force is considered to be impulsive.}
\end{qoptions}
\qn{6} The figure shows a snapshot of a vibrating string at $t=0$. The particle $P$ is observed moving up with a velocity $20\sqrt3$ cm/s. The tangent at $P$ makes an angle $60^\circ$ with $x$-axis. Then ($\mu$ mass per unit length of string $=50$ gm/m)
\qfig{p2_phy_s1_q06_1}
\begin{qoptions}
\opt{A}{the wave is moving along negative $x$-axis}
\opt{B}{the wave is moving along positive $x$-axis}
\opt{C}{the equation of wave is given by $y=0.4\sin2\pi\left(5t-\dfrac x4+\dfrac1{12}\right)$}
\opt{D}{the total energy carried by the wave per cycle is $3.6\times10^{-6}$ J}
\end{qoptions}
\sect{Section-I(ii)}
\qn{7} Match the following.
\par\noindent\textbf{List-I}
\begin{qoptions}
\opt{P}{Linear combination of two SHMs}
\opt{Q}{$y=A\sin\omega_1t+A\sin(\omega_2t+\phi)$}
\opt{R}{Time period of a pendulum of infinite length}
\opt{S}{Maximum value of time period of an Oscillating pendulum}
\end{qoptions}
\par\noindent\textbf{List-II}
\begin{qoptions}
\opt{1}{$T=2\pi\sqrt{\dfrac Rg}$ and not $T=\infty$}
\opt{2}{SHM for equal frequencies and amplitude}
\opt{3}{Superposition may not be a SHM always}
\opt{4}{Amplitude will be $\sqrt2A$ for $\omega_1=\omega_2$ and phase difference of $\dfrac\pi2$}
\end{qoptions}
\begin{qoptions}
\opt{A}{$P\to2,3$; $Q\to2,3,4$; $R\to1$; $S\to1$}
\opt{B}{$P\to1,3$; $Q\to2,3$; $R\to1,4$; $S\to1$}
\opt{C}{$P\to1,3$; $Q\to2,3,4$; $R\to1$; $S\to2,3$}
\opt{D}{$P\to2,4$; $Q\to1,3$; $R\to2,3$; $S\to3,4$}
\end{qoptions}
\qn{8} A monoatomic ideal gas of 2 moles was found to be undergoing a thermodynamic polytropic process whose volume coefficient is varying with temperature as shown. ($R=\dfrac{25}{3}$ J/mol/k.)
\qfig{p2_phy_s2_q08_1}
\par\noindent\textbf{List-I}
\begin{qoptions}
\opt{P}{Work done by gas in joules for every 1K rise in temperature, if $\tan\theta=3$}
\opt{Q}{Internal energy rise of the gas in joules for every 1K rise in temperature, if $\tan\theta=2$}
\opt{R}{Heat absorbed by the gas in joules for every 1K rise in temperature, if $\tan\theta=1$}
\opt{S}{Value of $\theta$ in degrees for which the internal energy rise of the gas for every 1 K rise in temperature is 25 J}
\end{qoptions}
\par\noindent\textbf{List-II}
\begin{qoptions}
\opt{1}{$0$}
\opt{2}{$25$}
\opt{3}{$\dfrac{125}{3}$}
\opt{4}{$\dfrac{50}{9}$}
\end{qoptions}
\begin{qoptions}
\opt{A}{$P\to4$; $Q\to2$; $R\to3$; $S\to1,2,3,4$}
\opt{B}{$P\to2,3$; $Q\to1,3,4$; $R\to2$; $S\to3,4$}
\opt{C}{$P\to3$; $Q\to3,4$; $R\to2,3,4$; $S\to1,2$}
\opt{D}{$P\to2,4$; $Q\to1,2$; $R\to4$; $S\to1,3,4$}
\end{qoptions}
\qn{9} A Young's double slit appratus is shown in figure
\qfig{p2_phy_s2_q09_1}
\par\noindent\textbf{List--I}
\begin{qoptions}
\opt{P}{Source $S$ move upwards}
\opt{Q}{A thin glass slab is placed in front of slit $S_2$}
\opt{R}{Whole apparatus is immersed in water}
\opt{S}{Diameter of slit $S_2$ is doubled}
\end{qoptions}
\par\noindent\textbf{List--II}
\begin{qoptions}
\opt{1}{Central fringe moves upwards}
\opt{2}{Central fringe moves downward}
\opt{3}{Fringe width remains constant}
\opt{4}{Fringe width decreases}
\opt{5}{Intensity of central maxima increases}
\end{qoptions}
\begin{qoptions}
\opt{A}{$P\to1,2$; $Q\to3,4$; $R\to2,5$; $S\to4$}
\opt{B}{$P\to2,3$; $Q\to2,3$; $R\to4$; $S\to3,5$}
\opt{C}{$P\to2,3$; $Q\to3,4$; $R\to1,4$; $S\to1,2$}
\opt{D}{$P\to1,2$; $Q\to2,3$; $R\to2,4$; $S\to1,3$}
\end{qoptions}
\qn{10} In list--I a situation is depicted each of which is in vertical plane. The surfaces are frictionless. Match these with appropriate entries in list--II.
\par\noindent\textbf{List-I}
\begin{qoptions}
\opt{P}{Bead is threaded on a circular fixed wire and is projected from the lowest point\\ \qfigi{p2_phy_s2_q10_1}}
\opt{Q}{Block loosely fits inside the fixed small tube and is protected from the lowest point\\ \qfigi{p2_phy_s2_q10_2}}
\opt{R}{Block is projected horizontally from the lowest point of a smooth fixed cylinder\\ \qfigi{p2_phy_s2_q10_3}}
\opt{S}{Block is projected on a fixed hemisphere from angular position $\theta$\\ \qfigi{p2_phy_s2_q10_4}}
\end{qoptions}
\par\noindent\textbf{List-II}
\begin{qoptions}
\opt{1}{Normal force is zero. At the top most point of its trajectory}
\opt{2}{Velocity of the body is zero. At top most point of it's trajectory}
\opt{3}{Acceleration of the body is zero at the top most point of it's trajectory}
\opt{4}{Normal force is radially outwards at the top most point of trajectory}
\end{qoptions}
\begin{qoptions}
\opt{A}{$P\to1,2$; $Q\to3,4$; $R\to2,5$; $S\to1,3$}
\opt{B}{$P\to2,3$; $Q\to2,3,4$; $R\to3,4$; $S\to1,2$}
\opt{C}{$P\to2,3$; $Q\to3,4$; $R\to2,3,4$; $S\to1,2$}
\opt{D}{$P\to2,4$; $Q\to1,2$; $R\to1$; $S\to2,3,4$}
\end{qoptions}
\sect{Section-II}
\qn{1} A large solenoid coil is made of single layer of conductor of cross section $4$cm $\times$ $2$cm with a cooling water hole $2$cm $\times$ $1$cm in the conductor. The coil, which consists of 100 turns, has a diameter of 3 meters, and a length of 4 meters. At the two ends of coil are circular steel plates to make the field uniform and to return the magnetic flux through a steel cylindrical structure to the coil, as shown in the diagram. A magnetic field of 0.25 tesla is desired. Resistivity of conductor is $3\times10^{-8}$ ohm-meter. What power (in megawatt) must be supplied to provide the desired field ?
\qfig{p2_phy_s3_q01_1}
\qn{2} To determine the density of blood, drops of blood were added to a mixture of xylene with a density of $0.867\dfrac{\mathrm{gm}}{\mathrm{cm}^3}$, and bromobenzene with a density of $1.497\dfrac{\mathrm g}{\mathrm{cm}^3}$. The ratio of xylene and bromobenzene was being changed until the drops were floating. i.e. they did not rise nor sink. This happened when the mixture contained $72\%$ of xylene and $28\%$ of bromobenzene by volume. What was the density (in gm/cm$^3$) of blood?
\qn{3} A $68$g bungee cord has an equilibrium length of $1.30$m. The cord is stretched to a length of $1.80$m, then vibrated at $20$ Hz. This produces a standing wave with two antinodes. What is the spring constant (in N/m) of bungee cord ? Assume both ends fixed?
\qn{4} A rod with length $3\ell$ and insignificant mass is suspended in horizontal position using two fibers attached to ends of the rod. There are two masses of $M$ and $2M$ fixed to the rod, so their centres are placed in one third or two thirds of the rod's length. Find the ratio of tension force in uncut string after and before cutting of string close to $2M$ will be :
\qfig{p2_phy_s3_q04_1}
\qn{5} Wire of an original length of $10$ m is fixed at one end and is being stretched with force of $200$ N which elongates it by $4$mm. Find the difference in diameter as multiple of $10^{-7}$ during the stretching of the wire. Where, Young's modulus $Y=2\times10^{11}$ N-m$^2$ and poisson's ratio $=0.32$
\qn{6} A conductor of length $\ell=10$ cm is placed in E-W direction in a plane. Earth's horizontal magnetic field is $B=3$ Tesla. Find the amount of charge (in mC) which passes through it when it is found to jump to height 5 cm. Mass of rod is 2 gm and $g=10$ m/s$^2$.
\qn{7} An isoceles right angled triangular glass ($\mu=1.6$) prism has a cavity inside it in the shape of a thin convex lens whose both surface have radius of curvature equal to $20$ cm. The cavity has been filled with a transparent liquid of refractive index $2.4$. $S$ is a point object and there is an opaque sheet having a small hole such that the source and the hole both lie on the principal axis of the lens. The small hole in the opaque sheet is just to ensure that only paraxial rays are incident on the original system. However, the size of hole is large enough to neglect diffraction effects. At what perpendicular distance (in cm) from line $BC$, the observer $P$ will see the image?
\qfig{p2_phy_s3_q07_1}
\qn{8} A non-conducting sphere of radius $R=5$cm has it's center at origin $O$ of co-ordinate system shown in figure. It has a spherical cavity of radius $r=1$ cm, whose centre is at $(0,3\text{cm})$. Solid material of sphere has uniform positive charge density $\rho=\dfrac{10^{-6}}{\pi}\mathrm{coul.\,m^{-3}}$. Assume dielectric constant of material to be 1. The potential (in volt) at point $P$ $(4\text{cm},0)$ is
\qfig{p2_phy_s3_q08_1}
\partsec{2}{CHEMISTRY}
\sect{Section-I(i)}
\qn{1} In a chemical reaction, the compound $(X)$
\qfig{p2_chem_s1_q01_1}
4(p--tolyl)--3 butenoic acid\par
Can rather conveniently be converted to another compound $(Y)$ which is
\qfig{p2_chem_s1_q01_2}
7--methyl 1--napthol\par
Which of the following is/are suitable for this purpose.
\begin{qoptions}
\opt{A}{Trans isomer of $(X)$ can be treated with anhydrous $\mathrm{AlCl_3}$.}
\opt{B}{Cis isomer of $(X)$ can be treated with HF.}
\opt{C}{Trans isomer of $(X)$ can be treated with conc.$\mathrm{H_3PO_4}$.}
\opt{D}{Cis isomer of $(X)$ can be treated with NaOH.}
\end{qoptions}
\qn{2} \qfigi{p2_chem_s1_q01_3}$\ \xrightarrow{\mathrm{MnO_2}}\ \mathrm A\ \xrightarrow[\mathrm{(ii)\,H_3O^+}]{\mathrm{(i)\,1\,Eq.,PhMgBr/CuCl}}\ \mathrm B$
\begin{qoptions}
\opt{A}{B is \qfigi{p2_chem_s1_q02_1}}
\opt{B}{A is \qfigi{p2_chem_s1_q02_2}}
\opt{C}{B is \qfigi{p2_chem_s1_q02_3}}
\opt{D}{A is \qfigi{p2_chem_s1_q02_4}}
\end{qoptions}
\qn{3} Which of the following will be soluble in $\mathrm{NaHCO_3}$ solution and more acidic than \qfigi{p2_chem_s1_q03_1}
\begin{qoptions}
\opt{A}{\qfigi{p2_chem_s1_q03_A}}
\opt{B}{\qfigi{p2_chem_s1_q03_B}}
\opt{C}{\qfigi{p2_chem_s1_q03_C}}
\opt{D}{\qfigi{p2_chem_s1_q03_D}}
\end{qoptions}
\qn{4} Which statements is/are correct among the following ?
\begin{qoptions}
\opt{A}{$\mathrm{BF_3}$ is strongest Lewis acid among $\mathrm{BF_3}$, $\mathrm{BCl_3}$, $\mathrm{BBr_3}$}
\opt{B}{Between $\mathrm{NH_3}$ and $\mathrm{PH_3}$, $\mathrm{NH_3}$ is better electron donor because $\mathrm{NH_3}$ has more directional lone pair.}
\opt{C}{In $\mathrm{P_4O_{10}}$ there are 5, P--O--P bonds.}
\opt{D}{$\mathrm{ClOF_3}$ has same shape as $\mathrm{XeO_2F_2}$.}
\end{qoptions}
\qn{5} Compound $(X)$, molecular formula $\mathrm{C_8H_8O}$, is water insoluble and gives yellow ppt. with $\mathrm{I_2}$/NaOH. It reacts with $\mathrm{H_2N-OH}$ in the presence of $\mathrm{NH_4Cl/NH_4OH}$ to give two geometrical isomers, which can go for rearrangement with acids. Compound $(X)$ can be,
\begin{qoptions}
\opt{A}{$\mathrm{C_6H_5{-}\overset{\overset{O}{\|}}{C}{-}CH_3}$}
\opt{B}{$\mathrm{C_6H_5{-}CH_2{-}CHO}$}
\opt{C}{$\mathrm{C_6H_5{-}\underset{OH}{CH}{-}CH_3}$}
\opt{D}{\qfigi{p2_chem_s1_q05_D}}
\end{qoptions}
\qn{6} In order to study effect of temperature on rate of reaction, another term ``$\sigma$'' is defined as $\sigma=\dfrac1k\dfrac{dk}{dT}$ [where $k$ is rate constant and $T$ represents temperature]. If for a reaction, the value of $\sigma=\dfrac{1.25\times10^6}{T^3}\mathrm{K^{-1}}$ [$T$ should be in Kelvin], then identify the correct statement(s).
\begin{qoptions}
\opt{A}{Activation energy for the given reaction at 250 Kelvin is 10 Kcal/mol.}
\opt{B}{Activation energy for the given reaction at 2000 Kelvin is 1.25 Kcal/mol.}
\opt{C}{Activation energy of this reaction increases on increasing temperature.}
\opt{D}{The value of $\dfrac{dk}{dT}$ at 2000 Kelvin is $0.625\times k$.}
\end{qoptions}
\sect{Section-I(ii)}
\qn{7} Match the following.
\par\noindent\textbf{List - I (Conversion)}
\begin{qoptions}
\opt{P}{$\mathrm{R{-}\overset{\overset{O}{\|}}{C}{-}NH_2\longrightarrow RNH_2}$}
\opt{Q}{$\mathrm{R{-}\overset{\overset{O}{\|}}{C}{-}NH{-}O{-}\overset{\overset{O}{\|}}{C}{-}R\longrightarrow RNH_2}$}
\opt{R}{$\mathrm{R{-}\overset{\overset{O}{\|}}{C}{-}OH\longrightarrow RNH_2}$}
\opt{S}{$\mathrm{R{-}\overset{\overset{O}{\|}}{C}{-}NH{-}NH_2\longrightarrow RNH_2}$}
\end{qoptions}
\par\noindent\textbf{List - II (Reagent)}
\begin{qoptions}
\opt{1}{KOBr}
\opt{2}{$\mathrm{OH^-/H_2O}$}
\opt{3}{$\mathrm{N_3H/H^+}$}
\opt{4}{$\mathrm{HNO_2/\Delta/H_3O^+}$}
\end{qoptions}
\par\noindent\textbf{Codes :}
\begin{qoptions}
\opt{A}{$P\to(1)$ \quad $Q\to(3)$ \quad $R\to(2)$ \quad $S\to(4)$}
\opt{B}{$P\to(1)$ \quad $Q\to(2)$ \quad $R\to(3)$ \quad $S\to(1)$}
\opt{C}{$P\to(1)$ \quad $Q\to(2)$ \quad $R\to(3)$ \quad $S\to(4)$}
\opt{D}{$P\to(2)$ \quad $Q\to(3)$ \quad $R\to(4)$ \quad $S\to(1)$}
\end{qoptions}
\qn{8} Match the following.
\par\noindent\textbf{List - I (Organic compound)}
\begin{qoptions}
\opt{P}{\qfigi{p2_chem_s2_q08_1}}
\opt{Q}{\qfigi{p2_chem_s2_q08_2}}
\opt{R}{\qfigi{p2_chem_s2_q08_3}}
\opt{S}{$\mathrm{C_2H_5NH_2}$}
\end{qoptions}
\par\noindent\textbf{List - II (Estimation/Detection Possible using)}
\begin{qoptions}
\opt{1}{Dumas method}
\opt{2}{Kjeldahl's method}
\opt{3}{Carius method}
\opt{4}{$p$--$\mathrm{CH_3C_6H_4SO_2Cl}$ (Hinsberg method)}
\end{qoptions}
\par\noindent\textbf{Codes :}
\begin{qoptions}
\opt{A}{$P\to(1,2,3)$ \quad $Q\to(3,4)$ \quad $R\to(1,2)$ \quad $S\to(1,2,4)$}
\opt{B}{$P\to(1,2,3,4)$ \quad $Q\to(3)$ \quad $R\to(1,2,3)$ \quad $S\to(1,2,4)$}
\opt{C}{$P\to(1,2,3)$ \quad $Q\to(1,4)$ \quad $R\to(2,3)$ \quad $S\to(4)$}
\opt{D}{$P\to(1,4)$ \quad $Q\to(3,4)$ \quad $R\to(2,3)$ \quad $S\to(1,3,4)$}
\end{qoptions}
\qn{9} Match the following. $A=$ Central atom, $X=$ Neighbouring atom(s)
\par\noindent\textbf{List - I}
\begin{qoptions}
\opt{P}{$\mathrm{PCl_3F_2}$ \& $\mathrm{SOF_4}$}
\opt{Q}{$\mathrm{BF_3}$ \& $\mathrm{BH_3}$}
\opt{R}{$\mathrm{NO_3^-}$ \& $\mathrm{NF_3}$}
\opt{S}{$\mathrm{XeF_5^+}$ \& $\mathrm{ICl_2^+}$}
\end{qoptions}
\par\noindent\textbf{List - II}
\begin{qoptions}
\opt{1}{$\angle X{-}A{-}X=120^\circ$ observed in both species}
\opt{2}{$\angle X{-}A{-}X$ less than $120^\circ$ observed in both species}
\opt{3}{Lone pair(s) at `$A$' in both species}
\opt{4}{Both species non--polar}
\opt{5}{Same `$A$' but different hybridization in both species}
\end{qoptions}
\par\noindent\textbf{Codes :}
\begin{qoptions}
\opt{A}{$P\to(1)$ \quad $Q\to(4,5)$ \quad $R\to(5)$ \quad $S\to(3)$}
\opt{B}{$P\to(1,2)$ \quad $Q\to(1,4)$ \quad $R\to(5)$ \quad $S\to(3)$}
\opt{C}{$P\to(2)$ \quad $Q\to(1,4,5)$ \quad $R\to(5)$ \quad $S\to(2,3)$}
\opt{D}{$P\to(2)$ \quad $Q\to(1,4)$ \quad $R\to(5)$ \quad $S\to(2,3)$}
\end{qoptions}
\qn{10} Match the following.
\par\noindent\textbf{List - I} (Electronic transition of hydrogenic ``sample'')
\begin{qoptions}
\opt{P}{$n=4$ to $n=2$ for $\mathrm{He^+}$}
\opt{Q}{$n=5$ to $n=2$ for $\mathrm H$}
\opt{R}{$n=6$ to $n=3$ for $\mathrm{Li^{2+}}$}
\opt{S}{$n=8$ to $n=4$ for $\mathrm H$}
\end{qoptions}
\par\noindent\textbf{List - II} (Observation for transition(s))
\begin{qoptions}
\opt{1}{Total 10 emission lines}
\opt{2}{``Balmer'' line(s) in visible spectrum}
\opt{3}{``Balmer'' line(s) in U--V spectrum}
\opt{4}{Total 6 emission lines}
\opt{5}{No ``Paschen'' line(s) observed}
\end{qoptions}
\par\noindent\textbf{Codes :}
\begin{qoptions}
\opt{A}{$P\to(2)$ \quad $Q\to(2,4)$ \quad $R\to(4)$ \quad $S\to(5)$}
\opt{B}{$P\to(3)$ \quad $Q\to(2)$ \quad $R\to(4)$ \quad $S\to(4,5)$}
\opt{C}{$P\to(3)$ \quad $Q\to(2,4)$ \quad $R\to(4)$ \quad $S\to(1,5)$}
\opt{D}{$P\to(2)$ \quad $Q\to(2,4)$ \quad $R\to(4,5)$ \quad $S\to(1)$}
\end{qoptions}
\sect{Section-II}
\qn{1} A $10.0$ gram sample of a mixture of calcium chloride and sodium chloride is treated with $\mathrm{Na_2CO_3}$. The precipitate formed is decomposed on heating. The weight of the residue is $1.68$ gram. The percentage by mass of $\mathrm{CaCl_2}$ in the original mixture is :
\qn{2} The emf of the following cell is given as $E_{\mathrm{Cell}}=-a\times10^{-1}$ V\\[2pt]
$\mathrm{Pt,H_2(1\ atm)}\,|\,$ $\mathrm{CH_3COOH(0.1\ M)}\,||\,$ $\mathrm{NH_4OH(0.1\ M)}\,|\,$ $\mathrm{H_2\ (1\ atm),\ Pt}$\\
$K_a$ for $\mathrm{CH_3COOH}=10^{-5}$; $K_b$ for $\mathrm{NH_4OH}=10^{-5}$, Given $\dfrac{2.303RT}{F}=0.06$ V\\
What is the value of $a$ ?
\qn{3} Consider the following two first order reaction :\\
$A\to P$ \hfill \ldots(i)\\
$B\to Q$ \hfill \ldots(ii)\\
Reaction (i) is $75\%$ complete in $4$ hrs. while reaction (ii) takes $16$ hrs. for same $75\%$ completion of reaction under identical conditions. By how many times, half life of (ii) is greater than the half life of (i) ?
\qn{4} A mixture contains $\mathrm{AgCl}$, $\mathrm{Al(OH)_3}$, $\mathrm{Zn(OH)_2}$, $\mathrm{Cu(OH)_2}$, $\mathrm{Fe(OH)_3}$ \& $\mathrm{Cd(OH)_2}$. On adding excess of $\mathrm{NH_4OH}$ how many of the metal complex cations will transfer into filtrate ? (Assume original mixture do not react among themselves)
\qn{5} Sum of tags of Nitrogen lone pairs which are delocaized.
\qfig{p2_chem_s3_q05_1}
\qn{6} Identify the species in which the ``actual'' as well as the ``average'' Oxidation number of the marked element is/are an \underline{EVEN integer}
\begin{center}
\begin{tabular}{@{}llll@{}}
(1) $\mathrm{H_4\underline{P}_2O_6}$ & (2) $\mathrm{Na_2\,\underline{S}_4O_6}$ & (3) $\mathrm{H_2\underline{S}\,O_5}$ & (4) $\mathrm{H_2\,\underline{S}_2O_8}$\\
(5) $\mathrm{H_2\underline{S}_2O_7}$ & (6) $\mathrm{\underline{S}_8}$ & (7) $\mathrm{\underline{Cr}O_8^{3-}}$ & (8) $\mathrm{H\underline{N}O_4}$\\
(9) $\mathrm{\underline{N}_2O_2^{2-}}$ & (10) $\mathrm{\underline{Os}O_4}$ & (11) $\mathrm{\underline{P}H_3}$ & (12) $\mathrm{\underline{Cr}O_4^{2-}}$\\
\end{tabular}
\end{center}
\qn{7} A solution is created by mixing equal volumes of $0.1$ M $\mathrm{Na_2SO_4}$ and $0.1$ M $\mathrm{BaCl_2}$. The resultant solution volume is $1$ litre, its temperature rises to `$T$' $=50^\circ$C and it also contains a precipitate mass. Through highly accurate measurements, total ionic concentration is found to be $0.20005$ M (excluding $\mathrm{H^+}$ and $\mathrm{OH^-}$ ions). Then report the number of correct statement(s) among the following
\begin{qoptions}
\opt{A}{Final solution contains equal moles of all ions (excluding $\mathrm{H^+}$ and $\mathrm{OH^-}$ ions)}
\opt{B}{Final solution contains equal moles of $\mathrm{Na^+}$ and $\mathrm{Cl^-}$ ions.}
\opt{C}{Final solution contains equal moles of $\mathrm{Ba^{+2}}$ and $\mathrm{SO_4^{2-}}$ ions.}
\opt{D}{Ksp value of $\mathrm{BaSO_4}$ is $2.5\times10^{-9}$ at temperature `$T$'}
\opt{E}{pH of resultant solution is 7}
\opt{F}{Resultant solution after filteration can act as buffer solution}
\opt{G}{Final solution is neutral because has $[\mathrm H^+]_{\mathrm{eqm}}=[\mathrm{OH^-}]_{\mathrm{eqm}}$}
\end{qoptions}
\qn{8} $40$ ml of $0.05$ M solution of sodium sesquicarbonate ($\mathrm{Na_2CO_3.NaHCO_3.\,2H_2O}$) is titrated against $0.05$ M HCl. ``$X$'' ml of HCl is used when phenolphthalein is the indicator and ``$Y$'' ml of HCl is used when methyl orange is the indicator in two separate titrations. Hence, value of $\dfrac{(Y-X)}{10}$ is
\partsec{3}{MATHEMATICS}
\sect{Section-I(i)}
\qn{1} If $a\ell^2+bm^2+2d\ell+1=0$, where $a,b,d$ are fixed real numbers such that $a-b=d^2$, then the line $\ell x+my+1=0$ touches the fixed circle (where $b<0$ and $\ell,m\in\mathbb R$)
\begin{qoptions}
\opt{A}{which cut the $x$-axis orthogonally}
\opt{B}{with radius equal to $\sqrt{-b}$}
\opt{C}{on which length of the tangent from origin is $\sqrt{d^2+b}$}
\opt{D}{with centre $\equiv(d,-b)$}
\end{qoptions}
\qn{2} $p(x)=x^2+px+q$ is a polynomial with complex coefficients such that $p(p(1))=p(p(2))=0$ and $p(1)\neq p(2)$. Then
\begin{qoptions}
\opt{A}{$|p(0)|=\dfrac12$}
\opt{B}{$p(0)=-\dfrac32$}
\opt{C}{$|p(1)|=1$}
\opt{D}{$p(x)=0$ has real roots}
\end{qoptions}
\qn{3} Let $f$ be a differentiable function such that $f(x)-f(y)=(x-y)f'\!\left(\dfrac{x+y}{2}\right)\ \forall\ x,y$ and $f(0)=f(1)=0$, $f(2)>0$. Then
\begin{qoptions}
\opt{A}{$f$ is onto}
\opt{B}{$f$ has extrema at $x=\dfrac12$}
\opt{C}{$f$ is many one}
\opt{D}{$\displaystyle\lim_{x\to-\infty}f(x)=-\infty$}
\end{qoptions}
\qn{4} Let $A,B$ be square matrices of order $n$. Then
\begin{qoptions}
\opt{A}{If $A,B,A+B$ are idempotent then $AB=BA=O$}
\opt{B}{If $A,B$ are orthogonal then $AB,BA$ are orthogonal}
\opt{C}{If $A+B=I$ and $A$ is idempotent then $AB=BA=O$}
\opt{D}{If $A^2+B^2=AB$ then $\operatorname{tr}(ABAB)=\operatorname{tr}(A^2B^2)$}
\end{qoptions}
\qn{5} Let $P=\left|z^5+\bar z^3+6z\right|-2\left|z^4+1\right|$ where $|z|=1$ and $Q$ be the number of real roots of $5-x=\sqrt{x+\sqrt{x+\sqrt x}}$, $z\in\mathbb C$, $x\in\mathbb R$. Which of the following are \textbf{FALSE}?
\begin{qoptions}
\opt{A}{$P_{\min}+Q=5$}
\opt{B}{$P_{\max}+Q=4$}
\opt{C}{$P+Q\in[4,5]$}
\opt{D}{$P-Q\in[3,4]$}
\end{qoptions}
\qn{6} Identify the correct statements.
\begin{qoptions}
\opt{A}{$g(x)+g(x+2)=10\ \forall x$ then $\displaystyle\int_{-1}^{7}g(x)\,dx=40$}
\opt{B}{$\displaystyle\lim_{\theta\to0^+}\int_0^\infty\dfrac{\tan^{-1}x}{x^2+2x\cos\theta+1}\,dx=\dfrac\pi4$}
\opt{C}{$\left|\displaystyle\prod_{r=1}^{n-1}(i-\alpha_r)\right|$ where $\alpha_r$ are imaginary $n$th roots of unity, can be $\sqrt2$}
\opt{D}{$\begin{vmatrix}a^2&b\sin A&c\sin A\\ b\sin A&1&\cos A\\ c\sin A&\cos A&1\end{vmatrix}=1$}
\end{qoptions}
\sect{Section-I(ii)}
\qn{7} Match the following.
\par\noindent\textbf{List-I}
\begin{qoptions}
\opt{P}{$\displaystyle\lim_{n\to\infty}\left(n^{-\frac12\left(1+\frac1n\right)}\left(1^1\,2^2\,3^3\cdots n^n\right)^{1/n^2}\right)=e^{-p/q}$, then $p/q$ is greater than}
\opt{Q}{$(7!)!$ is divisible by $(7!)^{k!}\cdot(6!)!$, then maximum value of $k$ is}
\opt{R}{$\displaystyle\lim_{x\to0^+}x^{(x^x-1)}\ \le$}
\opt{S}{If $\alpha,\beta,\gamma$ are roots of $x^3-2x^2-1=0$ and $T_n=\alpha^n+\beta^n+\gamma^n$, then $\dfrac{T_{11}-T_8}{T_{10}}$ is greater than}
\end{qoptions}
\par\noindent\textbf{List-II}
\begin{qoptions}
\opt{1}{$5$}
\opt{2}{$6$}
\opt{3}{$\dfrac15$}
\opt{4}{$\dfrac16$}
\opt{5}{$\dfrac14$}
\end{qoptions}
\par\noindent Codes:
\begin{qoptions}
\opt{A}{P$\to$3,4; Q$\to$2,4; R$\to$1; S$\to$4,5}
\opt{B}{P$\to$1,4; Q$\to$2; R$\to$1; S$\to$3,4}
\opt{C}{P$\to$3; Q$\to$2,4; R$\to$4; S$\to$1,5}
\opt{D}{P$\to$3,4; Q$\to$2; R$\to$1,2; S$\to$3,4,5}
\end{qoptions}
\qn{8} Match the following.
\par\noindent\textbf{List-I}
\begin{qoptions}
\opt{P}{If exhaustive set of values of $x$ satisfying $[5x]+[4x]+[3x]=2$ is $[a,b)$ then the integer in $[2^3a,3^3b)$ is ($[\cdot]$ denotes GIF)}
\opt{Q}{Number of critical points of $f(x)=\max\{x^3-1,x^2-1\}$ is $\le$}
\opt{R}{If $f(x)=\begin{cases}\dfrac{\{x\}-\{-x\}}{\{x\}+\{-x\}}, & x\notin\mathbb Z\\ \cos\pi x, & x\in\mathbb Z\end{cases}$ then number of points of discontinuity in $[-2,2]$ is}
\opt{S}{The value of $x$ for which $(\sin^{-1}x)^2-(\cos^{-1}x)^2\ge\dfrac{\pi^2}{4}$ is always $\le$}
\end{qoptions}
\par\noindent\textbf{List-II}
\begin{qoptions}
\opt{1}{$3$}
\opt{2}{$5$}
\opt{3}{$8$}
\opt{4}{$1$}
\opt{5}{$4$}
\end{qoptions}
\par\noindent Codes:
\begin{qoptions}
\opt{A}{P$\to$1,2,3; Q$\to$2,4; R$\to$1; S$\to$4,5}
\opt{B}{P$\to$1,2,3,5; Q$\to$1,2,3,5; R$\to$5; S$\to$1,2,3,4,5}
\opt{C}{P$\to$1,2,3,5; Q$\to$1,2,3; R$\to$5; S$\to$1,2,3,5}
\opt{D}{P$\to$1,3,5; Q$\to$1,2,3; R$\to$1; S$\to$1,2,3,4}
\end{qoptions}
\qn{9} Match the following.
\par\noindent\textbf{List-I}
\begin{qoptions}
\opt{P}{If $p,q$ are real numbers such that $p^2+q>q^2+p$, then $p^2+q$ is greater than}
\opt{Q}{If $2t^3-9t^2+30-a=0$ has three real distinct roots then possible integral value of $a$ is}
\opt{R}{If $f:[3,\infty)\to[1,\infty)$, $f(x)=\pi^{x(x-3)}$, then number of solutions of $f(x)=f^{-1}(x)$ is always greater than}
\opt{S}{$\displaystyle\sum_{r=1}^{\infty}\sin^{-1}\left[\dfrac{2r+1}{r(r+1)\left(\sqrt{r^2+2r}+\sqrt{r^2-1}\right)}\right]$ is always less than}
\end{qoptions}
\par\noindent\textbf{List-II}
\begin{qoptions}
\opt{1}{$3$}
\opt{2}{$5$}
\opt{3}{$-\dfrac12$}
\opt{4}{$-1$}
\opt{5}{$8$}
\end{qoptions}
\par\noindent Codes:
\begin{qoptions}
\opt{A}{P$\to$1,2,3; Q$\to$2,4; R$\to$1; S$\to$4,5}
\opt{B}{P$\to$1,2,3; Q$\to$1,3; R$\to$5; S$\to$2,5}
\opt{C}{P$\to$1,2,3,4; Q$\to$1,2,3; R$\to$5; S$\to$1,2,3,5}
\opt{D}{P$\to$3,4; Q$\to$2,5; R$\to$3,4; S$\to$1,2,5}
\end{qoptions}
\qn{10} Match the following.
\par\noindent\textbf{List-I}
\begin{qoptions}
\opt{P}{In $\triangle ABC$, $I$ is the incentre. $AI+BI+CI\ge\lambda r$, where $\lambda$ may be ($AI,BI,CI$ are distances from vertices and $r$ is the inradius)}
\opt{Q}{$f(x)=30-2x-x^3$ and $f_n=f\circ f\circ\cdots\circ f$ ($n$ times). Positive integral $x$ satisfying $f_{2020}(x)<f_{2019}(-x)$ is}
\opt{R}{Power of $x^2+xy+y^2$ dividing $(x+y)^7-x^7-y^7$ is}
\opt{S}{$a,b,c$ are complex numbers with $ab+bc+ca=\dfrac12$, $a+b+c=2$, $abc=4$, then $\left|9\left(\dfrac1{ab+c-1}+\dfrac1{ac+b-1}+\dfrac1{bc+a-1}\right)\right|$ is always less than}
\end{qoptions}
\par\noindent\textbf{List-II}
\begin{qoptions}
\opt{1}{$2$}
\opt{2}{$3$}
\opt{3}{$5$}
\opt{4}{$6$}
\opt{5}{$1$}
\end{qoptions}
\par\noindent Codes:
\begin{qoptions}
\opt{A}{P$\to$1,2,3; Q$\to$2,4; R$\to$1; S$\to$4,5}
\opt{B}{P$\to$1,2,3,4,5; Q$\to$1,5; R$\to$1,5; S$\to$2,3,4}
\opt{C}{P$\to$1,2,3,4; Q$\to$1,2,3; R$\to$5; S$\to$1,2,3,5}
\opt{D}{P$\to$3,4; Q$\to$2,5; R$\to$3,4; S$\to$1,2,5}
\end{qoptions}
\sect{Section-II}
\qn{1} $a_1=a_2=2$ and $\dfrac{2a_{n-1}a_n}{a_{n-1}a_{n+1}-a_n^2}=n^3-n$ for $n=2,3,4,\dots$ Find $\dfrac{[S]}{16}$ where $S=\displaystyle\sum_{k=2}^{2020}\dfrac{a_{k+1}}{a_k}$ and $[\cdot]$ denotes GIF.
\qn{2} $f(x)=x^2+g'(0)x+g''(1)$ and $g(y)=y^2+f''(y)+f'(y)-2$. Then $f(3)+\dfrac1{g(2)}$ equals
\qn{3} $S_n=\dfrac1{1^n}+\dfrac1{3^n}+\dfrac1{5^n}+\cdots\infty$ and $\displaystyle\int_0^1\dfrac{x(\ln x)^2}{1-x^4}\,dx=mS_3$. Find $m$.
\qn{4} $M=3a7b$ and $N=a37b$ are natural numbers where $a,b\in\{0,1,2,\dots,9\}$ such that $\dfrac M8$ \textbf{OR} $\dfrac N{12}$ are integers. If the probability that $\dfrac M{12}$ \textbf{AND} $\dfrac N8$ are also integers is $p$, then $4p=$
\qn{5} $P(t^2,2t)$, $1\le t\le3$ is a point on the parabola $y^2=4x$. $PM$ is drawn perpendicular to the tangent at the vertex. From the midpoint $Q$ of $PM$, $QL$ is drawn perpendicular to the focal chord through $P$. The maximum length of $QL$ is
\qn{6} Let $z_0=\dfrac{1-i}{2}$. Then \[\lim_{n\to\infty}\left|(1+z_0)\left(1+z_0^2\right)\left(1+z_0^4\right)\cdots\left(1+z_0^{2^n}\right)\right|,\ n\in\mathbb N\] must be
\qn{7} The gamma function is defined with $\Gamma(x+1)=x\Gamma(x)$, $\Gamma(1)=1$, $\Gamma\!\left(\tfrac12\right)=\sqrt\pi$ and the beta function $B(x,y)=\displaystyle\int_0^1t^{x-1}(1-t)^{y-1}\,dt=\dfrac{\Gamma(x)\Gamma(y)}{\Gamma(x+y)}$. If $\displaystyle\int_0^{\pi/2}\sin^6x\cos^8x\,dx=\dfrac{5\pi}{2^b}$ then $b$ equals
\qn{8} The plane $ax+by+cz=1$ containing the line $\dfrac{x-3}{2}=\dfrac{y-1}{4}=\dfrac{z-2}{5}$ is rotated through an angle $\dfrac\pi2$ about this line so as to contain the point $(4,3,7)$. The value of $\left(2a+3b+\dfrac c8\right)$ equals
\end{multicols}
\end{document}
